
\documentclass[letterpaper,11pt]{article}
%%%%%%%%%%%%%%%%%%%%%%%%%%%%%%%%%%%%%%%%%%%%%%%%%%%%%%%%%%%%%%%%%%%%%%%%%%%%%%%%%%%%%%%%%%%%%%%%%%%%%%%%%%%%%%%%%%%%%%%%%%%%%%%%%%%%%%%%%%%%%%%%%%%%%%%%%%%%%%%%%%%%%%%%%%%%%%%%%%%%%%%%%%%%%%%%%%%%%%%%%%%%%%%%%%%%%%%%%%%%%%%%%%%%%%%%%%%%%%%%%%%%%%%%%%%%
\usepackage{graphicx}
\usepackage{amsmath,amsfonts,amssymb,amsthm,float}
\usepackage{hyperref}
\usepackage[utf8]{inputenc}
\usepackage[left=2cm, right=2cm, top=2cm, bottom=2cm]{geometry}

\setcounter{MaxMatrixCols}{10}
%TCIDATA{OutputFilter=LATEX.DLL}
%TCIDATA{Version=5.50.0.2960}
%TCIDATA{<META NAME="SaveForMode" CONTENT="1">}
%TCIDATA{BibliographyScheme=BibTeX}
%TCIDATA{LastRevised=Tuesday, February 13, 2024 15:00:45}
%TCIDATA{<META NAME="GraphicsSave" CONTENT="32">}

\input{tcilatex}
\renewcommand{\baselinestretch}{1.15}
\setlength{\parindent}{0pt}
\setlength{\parskip}{0.5\baselineskip}
\pretolerance=2000 \tolerance=3000
\renewcommand{\abstractname}{Resumen}

\begin{document}

\title{Nombre de la pr\'{a}ctica}
\author{Nombre completo del alumno (No. Control) \\
%EndAName
Departamento de Ingenier\'{\i}a El\'{e}ctrica y Electr\'{o}nica\\
Tecnol\'{o}gico Nacional de M\'{e}xico / Instituto Tecnol\'{o}gico de Tijuana%
}
\maketitle

\noindent \textbf{Palabras clave: }Palabra clave 1; Palabra clave 2; Palabra
clave 3; Palabra clave 4; Palabra clave 5.

\noindent El autor debe identificar palabras clave espec\'{\i}ficas que
representen con precisi\'{o}n la investigaci\'{o}n, se debe evitar usar t%
\'{e}rminos generales o vagos que podr\'{\i}an aplicarse a muchos temas
diferentes; se sugiere escribir de cuatro a siete palabras clave ordenadas
alfab\'{e}ticamente, separadas por punto y coma, y omitir palabras que ya est%
\'{a}n en el t\'{\i}tulo del documento.

\bigskip

\noindent Correo: \textbf{xxx.xxx@tectijuana.edu.mx}

\noindent \noindent Carrera: \textbf{Ingenier\'{\i}a Biom\'{e}dica}

\noindent Asignatura: \textbf{Modelado de Sistemas Fisiol\'{o}gicos}

\noindent Profesor: \href{https://biomath.xyz/}{\textbf{Dr. Paul Antonio
Valle Trujillo}} (paul.valle@tectijuana.edu.mx)

\section{Funci\'{o}n de transferencia}

\subsection{Ecuaciones principales}

\subsection{Transformada de Laplace}

\subsection{Procedimiento algebraico}

\subsection{Resultado}

\section{Estabilidad del sistema en lazo abierto}

\section{Modelo de ecuaciones integro-diferenciales}

\section{Error en estado estacionario}

\begin{equation*}
e\left( t\right) =\lim_{s\rightarrow 0}sR\left( s\right) \left[ 1-\dfrac{%
V_{s}\left( s\right) }{V_{e}\left( s\right) }\right] =
\end{equation*}

\section{C\'{a}lculo de componentes para el controlador PID}

\begin{eqnarray*}
k_{I} &=&\frac{1}{R_{e}C_{r}} \\
k_{P} &=&\frac{R_{r}}{R_{e}} \\
k_{D} &=&R_{r}C_{e}
\end{eqnarray*}

\begin{equation*}
C_{r}=
\end{equation*}

\begin{eqnarray*}
R_{e} &=&\frac{1}{k_{I}C_{r}}= \\
R_{r} &=&k_{P}R_{e}= \\
C_{e} &=&\frac{k_{D}}{R_{r}}=
\end{eqnarray*}

\end{document}
